\documentclass[12pt,a4paper]{article}
\usepackage[utf8]{inputenc}
\usepackage[croatian]{babel}
\usepackage{geometry}
\geometry{margin=2.5cm}
\usepackage{graphicx}
\usepackage{hyperref}
\usepackage{enumitem}
\usepackage{booktabs}
\usepackage{xcolor}
\usepackage{listings}
\usepackage{float}

\definecolor{codegreen}{rgb}{0,0.6,0}
\definecolor{codegray}{rgb}{0.5,0.5,0.5}
\definecolor{codepurple}{rgb}{0.58,0,0.82}
\definecolor{backcolour}{rgb}{0.95,0.95,0.92}

\lstdefinestyle{mystyle}{
    backgroundcolor=\color{backcolour},   
    commentstyle=\color{codegreen},
    keywordstyle=\color{magenta},
    numberstyle=\tiny\color{codegray},
    stringstyle=\color{codepurple},
    basicstyle=\ttfamily\footnotesize,
    breakatwhitespace=false,         
    breaklines=true,                 
    captionpos=b,                    
    keepspaces=true,                 
    numbers=left,                    
    numbersep=5pt,                  
    showspaces=false,                
    showstringspaces=false,
    showtabs=false,                  
    tabsize=2
}
\lstset{style=mystyle}

\title{\textbf{Automatizacija sigurnosnog testiranja\\pomoću autonomnih AI agenata}}
\author{
Grupa 7 -- Podgrupa 2\\[1em]
\textit{Elektrotehnički fakultet Sarajevo}\\
\textit{Osnove informacionih sistema}
}
\date{Decembar 2025}

\begin{document}

\maketitle

\vspace{1cm}
\begin{center}
\textbf{Članovi Podgrupe 2:}\\[0.5em]
Enis Adilović, Danijal Alibegović, Emin Begić, Zana Beljuri, Bakir Činjarević
\end{center}

\newpage
\tableofcontents
\newpage

%============================================================
\section{Uvod}
%============================================================

\subsection{Cilj dokumenta}

Ovaj dokument istražuje primjenu autonomnih AI agenata u automatizaciji sigurnosnog testiranja softverskih sistema. Cilj je predstaviti kako moderne tehnike umjetne inteligencije mogu transformirati tradicionalne pristupe sigurnosnom testiranju, posebno u kontekstu kontinuirane integracije i isporuke (CI/CD).

\subsection{Kontekst i motivacija}

U današnjem softverskom razvoju, sigurnost više nije opcija već nužnost. Prema izvještaju IBM-a za 2023. godinu, prosječni trošak povrede podataka iznosi 4,45 miliona USD, dok se 83\% organizacija suočilo s više od jednog incidenta \cite{ibm2023}. Tradicionalno sigurnosno testiranje, koje se oslanja na manuelne preglede i periodične skenove, više nije dovoljno u okruženju gdje se novi kod deployira više puta dnevno.

\textbf{Ključni izazovi tradicionalnog pristupa:}
\begin{itemize}
    \item Manuelno testiranje ne skalira s brzinom modernog razvoja
    \item Sigurnosni stručnjaci su rijetki i skupi resurs
    \item Kasno otkrivanje ranjivosti povećava troškove popravke
    \item Statički alati generiraju previše lažnih pozitiva
\end{itemize}

Autonomni AI agenti nude rješenje kroz inteligentnu automatizaciju koja može kontinuirano analizirati kod, identificirati ranjivosti i čak predlagati popravke -- sve u realnom vremenu i bez ljudske intervencije.

%============================================================
\section{Pregled literature i tehnologija}
%============================================================

\subsection{Tradicionalne tehnike sigurnosnog testiranja}

Prije razmatranja AI pristupa, važno je razumjeti postojeće tehnike sigurnosnog testiranja:

\subsubsection{Static Application Security Testing (SAST)}

SAST analizira izvorni kod bez izvršavanja aplikacije. Alati poput \textbf{SonarQube}, \textbf{Checkmarx} i \textbf{Semgrep} skeniraju kod tražeći poznate obrasce ranjivosti (SQL injection, XSS, buffer overflow).

\textbf{Prednosti:} Rano otkrivanje, pokrivenost cijelog koda.\\
\textbf{Nedostaci:} Visoka stopa lažnih pozitiva, ne razumije runtime ponašanje.

\subsubsection{Dynamic Application Security Testing (DAST)}

DAST testira pokrenutu aplikaciju simulirajući napade izvana. Alati poput \textbf{OWASP ZAP}, \textbf{Burp Suite} i \textbf{Nikto} šalju maliciozne zahtjeve i analiziraju odgovore.

\textbf{Prednosti:} Pronalazi runtime ranjivosti, manje lažnih pozitiva.\\
\textbf{Nedostaci:} Zahtijeva pokrenutu aplikaciju, teško postiže punu pokrivenost.

\subsubsection{Interactive Application Security Testing (IAST)}

IAST kombinira SAST i DAST kroz instrumentaciju aplikacije. Alati poput \textbf{Contrast Security} prate tok podataka u realnom vremenu.

\subsubsection{Software Composition Analysis (SCA)}

SCA skenira zavisnosti (dependencies) za poznate ranjivosti. Alati poput \textbf{Snyk}, \textbf{Dependabot} i \textbf{OWASP Dependency-Check} provjeravaju biblioteke protiv CVE baza.

\subsubsection{Penetration Testing}

Manuelno ili polu-automatsko testiranje gdje etički hakeri pokušavaju probiti sigurnost sistema koristeći iste tehnike kao napadači.

\subsection{Uloga AI u sigurnosnom testiranju}

Primjena AI u sigurnosti nije nova, ali nedavni napredak u velikim jezičkim modelima (LLM) i autonomnim agentima otvara nove mogućnosti:

\begin{itemize}
    \item \textbf{ML za detekciju anomalija} -- Prepoznavanje neobičnih obrazaca u mrežnom saobraćaju ili ponašanju korisnika
    \item \textbf{NLP za analizu koda} -- Razumijevanje semantike koda, ne samo sintakse
    \item \textbf{Reinforcement Learning} -- Autonomno učenje strategija napada i odbrane
    \item \textbf{LLM-bazirani agenti} -- GPT-4, Claude i slični modeli koji mogu razumjeti kontekst, generirati testove i predlagati popravke
\end{itemize}

%============================================================
\section{Metod: AI agenti u CI/CD pipeline-u}
%============================================================

\subsection{Arhitektura integracije}

Autonomni AI agent za sigurnosno testiranje integrira se u CI/CD pipeline kao zaseban stadij (stage) koji se pokreće automatski pri svakom pushu koda.

\begin{figure}[H]
\centering
\fbox{\parbox{0.9\textwidth}{
\centering
\textbf{CI/CD Pipeline sa AI Security Agentom}\\[1em]
\texttt{Code Push} $\rightarrow$ \texttt{Build} $\rightarrow$ \texttt{Unit Tests} $\rightarrow$ \textbf{AI Security Scan} $\rightarrow$ \texttt{Deploy to Staging} $\rightarrow$ \textbf{AI DAST} $\rightarrow$ \texttt{Deploy to Prod}
}}
\caption{Pozicija AI sigurnosnih agenata u CI/CD toku}
\end{figure}

\subsection{Komponente AI sigurnosnog agenta}

Moderni AI agent za sigurnosno testiranje sastoji se od sljedećih komponenti:

\begin{enumerate}
    \item \textbf{Kodni analizator (Code Analyzer)}
    \begin{itemize}
        \item Prima diff promjena iz Git-a
        \item Koristi LLM za semantičku analizu koda
        \item Identificira potencijalne sigurnosne probleme u kontekstu
    \end{itemize}
    
    \item \textbf{Planer testova (Test Planner)}
    \begin{itemize}
        \item Na osnovu analize koda, AI generira plan testiranja
        \item Prioritizira testove prema procijenjenom riziku
        \item Adaptira strategiju na osnovu prethodnih nalaza
    \end{itemize}
    
    \item \textbf{Executor testova (Test Executor)}
    \begin{itemize}
        \item Pokreće generirane testove (fuzzing, injection pokušaji)
        \item Koristi tradicionalne alate (ZAP, Nuclei) orkestrirane od strane AI
        \item Dinamički prilagođava napade na osnovu odgovora
    \end{itemize}
    
    \item \textbf{Analizator rezultata (Result Analyzer)}
    \begin{itemize}
        \item Filtrira lažne pozitive koristeći kontekstualno razumijevanje
        \item Rangira nalaze prema stvarnom riziku
        \item Generiše izvještaje razumljive developerima
    \end{itemize}
    
    \item \textbf{Remediation Engine}
    \begin{itemize}
        \item Predlaže konkretne popravke koda
        \item Može automatski kreirati pull request sa fix-om
        \item Uči iz prihvaćenih/odbijenih prijedloga
    \end{itemize}
\end{enumerate}

\subsection{Primjer integracije: GitHub Actions}

\begin{lstlisting}[caption=Primjer GitHub Actions workflow-a sa AI sigurnosnim agentom]
name: Security Pipeline

on: [push, pull_request]

jobs:
  ai-security-scan:
    runs-on: ubuntu-latest
    steps:
      - uses: actions/checkout@v4
        with:
          fetch-depth: 0  # Za diff analizu
      
      - name: AI Code Analysis
        uses: ai-security-agent/scan@v1
        with:
          api-key: ${{ secrets.AI_AGENT_KEY }}
          mode: "code-review"
          severity-threshold: "medium"
      
      - name: Deploy to Staging
        if: success()
        run: ./deploy-staging.sh
      
      - name: AI DAST Scan
        uses: ai-security-agent/dast@v1
        with:
          target-url: "https://staging.example.com"
          scan-profile: "api-focused"
          
      - name: Fail on Critical
        if: steps.ai-dast.outputs.critical-count > 0
        run: exit 1
\end{lstlisting}

\subsection{Autonomno ponašanje agenta}

Ključna razlika između tradicionalnih alata i AI agenta je \textbf{autonomija}. Agent može:

\begin{itemize}
    \item \textbf{Samostalno odlučivati} koje dijelove koda detaljnije testirati
    \item \textbf{Učiti iz konteksta} -- razumije poslovnu logiku aplikacije
    \item \textbf{Adaptirati strategiju} -- ako početni napadi ne uspiju, pokušava alternativne pristupe
    \item \textbf{Komunicirati} -- objašnjava svoje nalaze na razumljiv način
    \item \textbf{Sarađivati} -- radi zajedno s developerima, ne protiv njih
\end{itemize}

%============================================================
\section{Glavni rezultati i primjeri}
%============================================================

\subsection{Postojeći alati i platforme}

Nekoliko kompanija već nudi AI-bazirane sigurnosne agente:

\begin{table}[H]
\centering
\begin{tabular}{lll}
\toprule
\textbf{Alat} & \textbf{Tip} & \textbf{Ključna karakteristika} \\
\midrule
Snyk DeepCode AI & SAST & LLM analiza za smanjenje lažnih pozitiva \\
GitHub Copilot Security & Code Review & Inline sigurnosne sugestije \\
Semgrep AI & SAST & Pravila generirana od AI \\
Pentera & Pen Testing & Autonomni penetration testing \\
Darktrace & Runtime & AI za detekciju anomalija \\
\bottomrule
\end{tabular}
\caption{Pregled AI sigurnosnih alata}
\end{table}

\subsection{Studija slučaja: Autonomni penetration testing}

Istraživanje na University of Illinois (2024) demonstriralo je GPT-4 agenta koji je uspješno eksploatisao 87\% poznatih CVE ranjivosti u testnom okruženju, bez prethodnog znanja o specifičnim ranjivostima \cite{fang2024}.

Agent je koristio sljedeći pristup:
\begin{enumerate}
    \item Izviđanje (reconnaissance) ciljnog sistema
    \item Analiza prikupljenih informacija
    \item Generiranje hipoteza o mogućim ranjivostima
    \item Kreiranje i izvršavanje eksploita
    \item Iteracija na osnovu rezultata
\end{enumerate}

\subsection{Praktični primjer: SQL Injection detekcija}

Tradicionalni SAST alat bi označio sljedeći kod kao potencijalno ranjiv:

\begin{lstlisting}[language=Python, caption=Potencijalno ranjiv kod]
def get_user(user_id):
    query = f"SELECT * FROM users WHERE id = {user_id}"
    return db.execute(query)
\end{lstlisting}

Međutim, alat ne zna da li je \texttt{user\_id} validiran prije poziva. AI agent može:

\begin{enumerate}
    \item Analizirati cijeli tok podataka (data flow)
    \item Provjeriti da li postoji validacija na API nivou
    \item Ako ne postoji, predložiti konkretnu popravku:
\end{enumerate}

\begin{lstlisting}[language=Python, caption=AI-predložena popravka]
def get_user(user_id: int):  # Type hint kao prva linija odbrane
    query = "SELECT * FROM users WHERE id = :id"
    return db.execute(query, {"id": user_id})  # Parametrizovani upit
\end{lstlisting}

\subsection{Metrike uspješnosti}

Prema dostupnim istraživanjima, AI agenti pokazuju značajna poboljšanja:

\begin{itemize}
    \item \textbf{Smanjenje lažnih pozitiva:} 40-60\% u odnosu na tradicionalne SAST alate
    \item \textbf{Povećanje pokrivenosti:} Pronalaze 15-25\% više stvarnih ranjivosti
    \item \textbf{Brzina:} Analiza pull requesta za 2-5 minuta vs. sati za manuelni pregled
    \item \textbf{Developer experience:} 78\% developera preferira AI sugestije nad tradicionalnim izvještajima
\end{itemize}

%============================================================
\section{Zaključak i preporuke}
%============================================================

\subsection{Zaključak}

Autonomni AI agenti predstavljaju značajan napredak u automatizaciji sigurnosnog testiranja. Njihova sposobnost da razumiju kontekst, uče iz iskustva i autonomno donose odluke čini ih idealnim za integraciju u moderne CI/CD pipeline-e.

Ključne prednosti uključuju:
\begin{itemize}
    \item Kontinuirano testiranje bez ljudske intervencije
    \item Smanjenje lažnih pozitiva kroz kontekstualno razumijevanje
    \item Automatsko generiranje popravki koda
    \item Skalabilnost s brzinom razvoja
\end{itemize}

Međutim, AI agenti nisu zamjena za ljudske stručnjake. Najbolji rezultati postižu se hibridnim pristupom gdje AI obavlja repetitivne zadatke, a sigurnosni stručnjaci fokusiraju se na kompleksne probleme i strateške odluke.

\subsection{Preporuke za implementaciju}

Za organizacije koje žele implementirati AI sigurnosne agente, preporučujemo:

\begin{enumerate}
    \item \textbf{Početi s malim} -- Integrisati AI agenta na jednom projektu prije širenja
    \item \textbf{Definisati metrike uspjeha} -- Pratiti stopu lažnih pozitiva, vrijeme do popravke, broj pronađenih ranjivosti
    \item \textbf{Obučiti tim} -- Developeri trebaju razumjeti kako koristiti AI preporuke
    \item \textbf{Iterativno poboljšavati} -- Koristiti feedback loop za fino podešavanje agenta
    \item \textbf{Zadržati ljudski nadzor} -- Kritične odluke i dalje trebaju ljudsku verifikaciju
\end{enumerate}

\subsection{Budući pravci}

Očekujemo da će se AI sigurnosni agenti razvijati u smjeru:
\begin{itemize}
    \item Potpuno autonomnog penetration testinga
    \item Proaktivne zaštite (predviđanje napada prije nego se dogode)
    \item Integracije s bug bounty programima
    \item Regulatorne usklađenosti (automatska provjera GDPR, PCI-DSS)
\end{itemize}

%============================================================
\section*{Reference}
%============================================================
\addcontentsline{toc}{section}{Reference}

\begin{thebibliography}{9}

\bibitem{ibm2023}
IBM Security (2023). \textit{Cost of a Data Breach Report 2023}. IBM Corporation.

\bibitem{fang2024}
Fang, R., Bindu, R., Gupta, A., \& Kang, D. (2024). LLM Agents can Autonomously Exploit One-day Vulnerabilities. \textit{arXiv preprint arXiv:2404.08144}.

\bibitem{owasp2023}
OWASP Foundation (2023). \textit{OWASP Top 10:2021}. \url{https://owasp.org/Top10/}

\bibitem{gartner2024}
Gartner (2024). \textit{Market Guide for Application Security Testing}. Gartner Research.

\bibitem{snyk2023}
Snyk (2023). \textit{State of Open Source Security Report 2023}. Snyk Ltd.

\bibitem{microsoft2024}
Microsoft Security (2024). \textit{AI-Powered Security: The Next Frontier}. Microsoft Corporation.

\bibitem{github2024}
GitHub (2024). \textit{The State of Octoverse: Security}. GitHub, Inc.

\end{thebibliography}

\end{document}

